\documentclass[12pt,a4paper]{article}

\usepackage[utf8]{inputenc}
\usepackage[T1]{fontenc}
\usepackage[margin=2.5cm]{geometry}
\usepackage{amsmath,amssymb}
\usepackage{graphicx}
\usepackage{booktabs}
\usepackage{hyperref}
\usepackage{xcolor}
\usepackage{enumitem}
\usepackage{array}
\usepackage{multirow}
\usepackage{float}
\usepackage{microtype}

\hypersetup{
    colorlinks=true,
    linkcolor=blue!70!black,
    urlcolor=blue!70!black,
    citecolor=blue!70!black
}

\title{\textbf{DiffuVQA: Diffusion-Based Medical Visual Question Answering}\\[0.4em]
\large{Ara Dönem İlerleme Raporu --- Mayıs 2026}}
\author{Ali Ekin Özçetin}
\date{19 Mayıs 2026}

\begin{document}
\maketitle

% ─────────────────────────────────────────────
\section{Mimari}
% ─────────────────────────────────────────────

\subsection{Genel Boru Hattı}

\begin{enumerate}
    \item \textbf{Görüntü Kodlayıcı (CLIP ViT-B/32):} Tıbbi görüntü 384px çözünürlükte $\mathbb{R}^{d_v}$ boyutunda görsel temsile dönüştürülür.
    \item \textbf{Soru Kodlayıcı (BERT-base-uncased):} Soru tokenlara ayrılıp $\mathbb{R}^{L_q \times d_h}$ boyutunda bağlamsal temsile kodlanır.
    \item \textbf{Çapraz Dikkat Füzyonu:} Görsel ve sözel özellikler çapraz dikkat mekanizmasıyla bir koşullama vektörüne $\mathbf{c}$ birleştirilir.
    \item \textbf{Difüzyon (TransformerNetModel):} Cevap embedding'i $\mathbf{x}_T \sim \mathcal{N}(0, I)$'den başlayarak koşullu ters süreçle üretilir.
    \item \textbf{Çözme (Nearest-Neighbor Rounding):} Denoised embedding $\hat{\mathbf{x}}_0$ en yakın kelime dağarcığı token'ına eşlenerek metin üretilir.
\end{enumerate}

\subsection{Temel Konfigürasyon}

\begin{table}[H]
\centering
\begin{tabular}{ll}
\toprule
\textbf{Parametre} & \textbf{Değer} \\
\midrule
Difüzyon adımı & 2000 \\
Gürültü takvimi & sqrt \\
Gizli boyut & 768 \\
Dizi uzunluğu (seq\_len) & 32 \\
Öğrenme oranı & $5 \times 10^{-5}$ \\
Batch boyutu & 64 (microbatch 16) \\
Hedef eğitim adımı & 200\,000 \\
Zamanlayıcı örnekleyici & lossaware \\
EMA oranı & 0.9999 \\
\bottomrule
\end{tabular}
\caption{Mevcut eğitim konfigürasyonu.}
\end{table}

% ─────────────────────────────────────────────
\section{Uygulanan Mimari Değişiklikler}
% ─────────────────────────────────────────────

Bu bölümde proje boyunca alınan ve uygulanan mimari kararlar önem sırasına göre özetlenmektedir.

\subsection{Kritik Hata Düzeltmeleri}

\subsubsection{lm\_head Ağırlık Bağlama (Weight Tying) Hatası}

\textbf{Sorun:} BERT dalında \texttt{word\_embedding} pretrained ağırlıkla güncellenince Python referans semantiği nedeniyle \texttt{lm\_head.weight} bağlantısı kopuyordu. \texttt{lm\_head} rastgele başlatılmış halde kalıyordu.

\textbf{Sonuç:} \texttt{get\_efficient\_knn} ve \texttt{get\_logits} fonksiyonları rastgele bir matrise karşı L2 mesafesi hesaplıyordu. 100K adım sonrasında bile $\text{avg\_nn\_l2} = 558$ --- embedding manifoldu hiç öğrenilemiyordu.

\textbf{Çözüm:} Her pretrained dal (\texttt{bert}, \texttt{pubmedbert}, \texttt{roberta}) için \texttt{word\_embedding} atandıktan \textit{sonra} \texttt{lm\_head.weight = word\_embedding.weight} açıkça tekrar set edildi.

\subsubsection{Microbatch Gradyan Biriktirme Hatası}

\textbf{Sorun:} \texttt{backward()}, \texttt{schedule\_sampler.update\_with\_local\_losses()} ve \texttt{log\_loss\_dict()} microbatch döngüsünün dışında çağrılıyordu; yalnızca son microbatch'in gradyanı geriye yayılıyordu.

\textbf{Sonuç:} batch=64 ve microbatch=16 ile 4 yerine 1 etkin backward çalışıyordu. Öğrenme kalitesi 4 kat düşmüş görünüyordu.

\textbf{Çözüm:} İlgili çağrılar döngü içine alındı; loss $/ N_{\text{micro}}$ ile ölçeklendi.

\subsubsection{Eğitim/Çıkarım Uyumsuzluğu (use\_noising\_f)}

\textbf{Sorun:} Eğitimde ileri süreç cevap embedding'ini ($\mathbf{f} = \mathbf{x}_{\text{ans}}$) görüyordu; çıkarımda bu bilgi yoktu.

\textbf{Çözüm:} \texttt{use\_noising\_f=False} bayrağı eklendi. Eğitim de çıkarım gibi pure noise'dan başlıyor.

\subsection{Loss Fonksiyonu İyileştirmeleri}

\subsubsection{Padding Mask ile Loss Maskeleme}

SLAKE cevaplarının ortalama uzunluğu 1.5 token; seq\_len=32 ile 30 pozisyon padding'dir.

\textbf{Önceki:} MSE loss tüm 32 pozisyon için hesaplanıyordu --- model padding pozisyonlarında anlamsız token üretmeyi öğreniyor, gerçek cevap pozisyonlarına odaklanamıyordu.

\textbf{Sonraki:}
\begin{equation}
\mathcal{L}_{\text{MSE}} = \frac{\sum_{i=1}^{L} m_i \cdot \|\mathbf{x}_0^{(i)} - \hat{\mathbf{x}}_0^{(i)}\|^2}{\sum_{i=1}^{L} m_i}
\end{equation}
burada $m_i = \mathbf{1}[\text{token}_i \neq \text{PAD}]$. NLL terimi de aynı maskeyle uygulandı.

\subsubsection{decoder\_nll ve tT\_loss Çıkarma}

\textbf{tT\_loss:} $\mathbb{E}[q(\mathbf{x}_T|\mathbf{x}_0)^2]$ --- ileri sürecin sonunda Gaussian'ın sıfıra yakınlığını ölçüyor, gereksiz kısıt koyuyor.

\textbf{decoder\_nll:} Temiz embedding'den hesaplanan NLL --- denoised çıktı üzerinden hesaplanan \texttt{terms["nll"]} ile çifte sayım yapıyordu.

\textbf{Güncel loss:}
\begin{equation}
\mathcal{L} = \mathcal{L}_{\text{MSE}} + \mathcal{L}_{\text{NLL}}
\end{equation}

\subsubsection{logits\_mode=2}

Çıkarım sırasında \texttt{get\_efficient\_knn} L2 metriği kullanıyor; eğitimde \texttt{logits\_mode=1} dot-product kullanıyordu. Tutarsızlık giderildi: \texttt{logits\_mode=2} ile eğitim ve çıkarım aynı metriği kullanır.

\subsection{Encoder Dondurma (Freeze) Kararları}

\subsubsection{CLIP Görüntü Kodlayıcısı (Donduruldu)}

\textbf{Gerekçe:} CLIP ViT-B/32 $\approx$151M parametredir. Tıbbi görüntüler için zengin görsel temsil zaten mevcut. 14K örnekte eğitim: overfitting riski ve 2GB ekstra bellek maliyeti.

\textbf{Uygulama:}
\begin{verbatim}
for p in self.vision_encoder.parameters():
    p.requires_grad_(False)
\end{verbatim}

\subsubsection{BERT Dil Kodlayıcısı (Donduruldu --- Yeni)}

\textbf{Sorun:} 30K adım sonrasında model tüm cevaplar için yalnızca ``the'' / ``in'' üretiyordu. $\text{avg\_nn\_l2}$: 559 (15K) $\to$ 562 (30K) --- iyileşme yok.

\textbf{Kök Neden Analizi:}

\begin{itemize}
    \item BERT-base-uncased 110M parametredir.
    \item CLIP + BERT + fusion + diffusion aynı anda optimize edilince model, yüksek frekanslı ``the''/``in'' tokenlarını üreterek loss'u minimize eden dejenere çözüme yakınsar.
    \item Padding mask uygulandı, seq\_len 64'ten 32'ye düşürüldü --- fakat bu tek başına yeterli değil. Collapse çıkarım sırasında oluyor: model anlamlı bir token dağılımı öğrenememiş.
\end{itemize}

\textbf{Neden seq\_len=16 Reddedildi:} Collapse eğitim loss'undan değil, token üretim dağılımından kaynaklanıyor. Dizi uzunluğunu düşürmek aynı dejenere çözümü daha kısa dizi üzerinden üretir.

\textbf{Uygulama:}
\begin{verbatim}
self.language_encoder = language_encoder
for p in self.language_encoder.parameters():
    p.requires_grad_(False)
\end{verbatim}

Yalnızca fusion + diffusion katmanları ($\approx$50M param) eğitilebilir olarak kaldı.

\subsection{Öğrenme Oranı ve EMA İyileştirmeleri}

\subsubsection{Warmup + Cosine Decay + Floor}

\begin{equation}
\text{lr}(t) = \begin{cases}
\text{lr}_{\text{base}} \cdot \frac{t}{T_{\text{warm}}} & t < T_{\text{warm}} \\[4pt]
\text{lr}_{\min} + \frac{1}{2}(\text{lr}_{\text{base}} - \text{lr}_{\min})\left(1 + \cos\!\left(\pi \cdot \frac{t - T_{\text{warm}}}{T - T_{\text{warm}}}\right)\right) & t \geq T_{\text{warm}}
\end{cases}
\end{equation}

$T_{\text{warm}} = 0.03 \times T$ (150K'da $\approx$4500 adım), $\text{lr}_{\min} = 0.05 \times \text{lr}_{\text{base}}$.

\subsubsection{Dinamik EMA Warmup}

\begin{equation}
r(t) = \min\!\left(r_{\text{target}},\; 1 - \frac{1}{t+1}\right)
\end{equation}

$t < 10\,000$ için geçerli. Step=1'de $r=0.5$ (50/50 mix), step=9999'da $r=0.9999$. Rastgele başlatılmış ağırlıkların EMA modeline gereğinden fazla katkısını önler.

\subsection{Çıktı Sonrası İşleme (Post-processing)}

\texttt{sample\_vqa\_GPU.py}'ye eklenenler:
\begin{itemize}
    \item \textbf{SEP/PAD kesme:} Üretilen dizide ilk \texttt{[SEP]} veya \texttt{[PAD]}'e kadar kes.
    \item \textbf{Güven eşiği:} Trailing token'larda güven $< 0.1$ olanları sil.
    \item \textbf{Top-k reranking} (\texttt{decode\_top\_k=5}): Her pozisyon için $k$ aday token üret, en yüksek skorlu diziyi seç.
    \item \textbf{MBR (Minimum Bayes Risk):} $N$ örnek üret, ortalama embedding'e L2 uzaklığı minimum olan seç.
\end{itemize}

% ─────────────────────────────────────────────
\section{Tartışılacak Konular}
% ─────────────────────────────────────────────

\subsection{Veri Seti Değişikliği}

\subsubsection{SLAKE'in Sınırlılıkları}

\begin{itemize}
    \item 14K örnek --- diffusion modeli için küçük.
    \item Ortalama cevap uzunluğu \textbf{1.5 token}: 32 pozisyonun 30'u padding. Model uzun, anlamlı cevap üretmeyi öğrenemiyor.
    \item Kapalı uçlu (\%Y/N) sorular ağırlıkta --- difüzyon tabanlı üretim bu kısıtlı uzayda avantaj sağlamıyor.
\end{itemize}

\subsubsection{Aday Veri Setleri}

\begin{table}[H]
\centering
\begin{tabular}{p{3.2cm}p{2cm}p{2cm}p{6cm}}
\toprule
\textbf{Veri Seti} & \textbf{Boyut} & \textbf{Ort. Cevap} & \textbf{Not} \\
\midrule
Kvasir-VQA-x1       & 6.5K  & 6--10 token & Endoskopi; uzun tanımlayıcı cevaplar. HuggingFace'te mevcut. Mimari değişiklik gerektirmiyor. \\
PathVQA             & 32.8K & 1/8 token   & Patoloji görüntüleri; SLAKE'e benzer kısa cevap. \\
VQA-RAD             & 3.5K  & 4--8 token  & Radyoloji; küçük ama kaliteli referans cevaplar. \\
Med-VQA 2019        & 12.8K & 2--5 token  & Genel tıbbi VQA; mevcut pipeline ile uyumlu. \\
PubMedVision        & 1.3M  & ---         & Büyük ölçekli; ön eğitim için uygun, fine-tune maliyetli. \\
\bottomrule
\end{tabular}
\caption{Aday veri setleri ve temel özellikleri.}
\end{table}

\subsubsection{Öneri: Kvasir-VQA-x1}

\begin{itemize}
    \item \textbf{Avantaj:} Uzun cevaplar difüzyon modelinin açık uçlu üretim kapasitesini gerçekten test eder.
    \item \textbf{Maliyet:} Görüntü formatı (endoskopi JPEG, 720p) CLIP ViT-B/32 ile uyumlu --- encoder değişikliği yok. HuggingFace'ten indirme ve \texttt{vqa\_datasets.py} için yeni bir JSONL dönüştürücü gerekiyor (1 günlük iş). A100'de eğitim süresi SLAKE ile benzer ($\approx$8--10 saat/100K adım).
    \item \textbf{Risk:} Daha az medikal etiket çeşitliliği --- genel Med-VQA benchmark'larında SLAKE kadar temsil edilmiyor.
\end{itemize}

\subsubsection{Değişikliğin Zamanlaması}

BERT freeze ile yeniden eğitim başladı. \textbf{Öneri:} 50K adımda örnekleme yapılıp $\text{avg\_nn\_l2} < 300$ ve üretilen cevaplarda çeşitlilik gözlenirse mevcut mimarinin doğru çalıştığı teyit edilir. Bu sonuç olumlu ise veri seti değişikliği bir sonraki run'a ertelenebilir; olumsuz ise hem mimari hem veri seti sorunun birbirini maskelediği değerlendirilmeli.

\subsection{Açık Teknik Sorular}

\begin{enumerate}
    \item \textbf{seq\_len seçimi:} SLAKE ort. cevap 1.5 token. Padding mask loss maskeleme uygulandı. Ancak çıkarım hâlâ seq\_len kadar token üretiyor. Üretim seq\_len'ini dinamik yapma (örn. ilk SEP'te dur) ne zaman uygulanmalı?
    \item \textbf{NLL ağırlığı:} 200K adımda NLL/MSE oranı izleniyor. $> 5\times$ olursa \texttt{nll\_weight} parametresi eklenmeli mi?
    \item \textbf{pre\_answer\_loss ablasyonu:} \texttt{pre\_answer\_loss\_weight=0.0} şu an. Fusion supervision (\texttt{0.05}) eğitim dinamiklerini iyileştirir mi?
    \item \textbf{Veri Seti:} Kvasir-VQA-x1'e geçiş ne zaman? BERT freeze sonuçları beklenecek mi, yoksa paralel run yapılacak mı?
\end{enumerate}

% ─────────────────────────────────────────────
\section*{Özet Tablo}
% ─────────────────────────────────────────────

\begin{table}[H]
\centering
\begin{tabular}{p{5.5cm}p{4cm}p{4cm}}
\toprule
\textbf{Değişiklik} & \textbf{Gerekçe} & \textbf{Durum} \\
\midrule
lm\_head weight tying restore  & avg\_nn\_l2 divergence fix & Uygulandı \\
Microbatch gradient fix         & 4x öğrenme kaybı   & Uygulandı \\
use\_noising\_f=False           & Eğitim/çıkarım hizalaması & Uygulandı \\
Padding mask loss maskeleme     & Kısa cevap optimizasyonu & Uygulandı \\
CLIP freeze                     & 151M param, overfitting & Uygulandı \\
BERT freeze                     & the/in collapse, 110M param & Uygulandı (yeni) \\
decoder\_nll + tT\_loss çıkarma & Çifte sayım, gereksiz kısıt & Uygulandı \\
logits\_mode=2                  & Eğitim/çıkarım metrik tutarlılığı & Uygulandı \\
Cosine LR + warmup + floor      & Erken stabilizasyon & Uygulandı \\
Dinamik EMA warmup              & Random init absorpsiyonu & Uygulandı \\
seq\_len: 64 $\to$ 32            & SLAKE kısa cevap yapısı & Uygulandı \\
seq\_len: 32 $\to$ 16            & Analiz: gereksiz             & Reddedildi \\
Kvasir-VQA veri seti            & Uzun cevap, 6.5K örnek       & Tartışılacak \\
\bottomrule
\end{tabular}
\caption{Tüm mimari kararların özeti.}
\end{table}

\end{document}
